% !TEX root = ../main.tex
\chapter{Explore: find a span}\label{ch:explore}
Explore is where you look at the recordings before you analyse anything. Its job is to help
you find a \term{span}{span}, a stretch of one \term{channel}{channel} worth a closer look, and
then to hand that span on: to Analyse, to run a chain over it, or to Review, to judge it.

The four pages go from wide to narrow. \apppage{explore.corpus} shows every channel of one
recording at once, coloured by how much has been marked where. \apppage{explore.signal} opens
one channel at its real scale, and is where you pick and send a span.
\apppage{explore.cross-channel} draws the same short window on several channels together,
with the time lag between them. \apppage{explore.span-edit} is the odd one out: you reach it
from Review, to redraw the edges of a span, and it still runs on example data.

\begin{plainterms}[Annotations and detections]
Two kinds of marks appear on every Explore page. An \term{annotation}{annotation} is a mark a
person made: a span with a \term{verdict}{verdict} such as \emph{interesting}. A
\term{detection}{detection} is a mark a machine made: a span a run of a detection chain
found. Think of a teacher's red pen and a spell-checker's underline on the same essay. The
app keeps them in separate tables and never mixes them (\cref{ch:review} explains why).
\end{plainterms}

\begin{plainterms}[Held-out data]
One recording, \code{M4_aug_concat_fs1.mat}, is \term{held-out}{held out}. It is kept
locked away for the final test of any model, like an exam paper nobody may see before the
exam. If you looked at it while building a method, you could tune the method to it without
meaning to, and the final test would no longer be fair. So every Explore page shows it as a
locked card, and the bridge refuses every request for its data.
\end{plainterms}

\pagesection{explore.corpus}{Corpus}{A bird's-eye view across every channel}{live}
Corpus draws one recording as a grid: a row per channel, a column per stretch of time. Each
cell is shaded by how many spans fall in it. You use it to see where people and machines
have been looking, and where they disagree.

\pagefield{When you'd come here}
At the start of a session, to choose a recording and a channel. Also when you want to know
which parts of a recording nobody has reviewed yet, or where a detection run and the human
labels disagree.

\screenshot{pages/explore/explore.corpus--default}{The coverage map of one recording. Darker
cells hold more spans; the outlined row is the selected channel, and the bottom bar repeats
its counts.}

\pagefield{Doing things here}
\begin{enumerate}
  \item \textbf{Choose a recording.} Click the recording selector at the top left. The menu
    lists every recording with its channel count, length and sampling rate. Type in
    \ui{filter recordings} to narrow it. A row marked \ui{held out · locked} cannot be
    picked. A row marked \ui{linked fs2} (or similar) is the same recording resampled at
    another rate, so its counts are not independent of its partner's. The arrows beside the
    selector page through the recordings one at a time.
  \item \textbf{Read the map.} Hover a cell to see the channel, the hours it covers and its
    count. The shading is by rank among the non-empty cells, not by raw count, so one very
    busy cell does not wash the rest out. Grey cells hold nothing. The \ui{bin} menu sets how
    many hours each column covers; \ui{auto} uses 57 columns.
  \item \textbf{Choose what is counted.} \ui{colour by} offers \ui{annotations},
    \ui{detections}, \ui{annotations + detections} (the sum of the two, not the cells where
    both occur) and \ui{disagree} (see below). In the right-hand rail, \ui{Show} turns the
    two layers on or off; \ui{Detections from} limits the detections to chosen runs and
    methods; \ui{Verdict} limits the annotations to chosen verdicts; \ui{Morphology tag}
    dims the channels that carry none of the chosen tags. The foot of the rail says how many
    spans pass every filter, on how many channels.
  \item \textbf{Look for unreviewed stretches.} Tick \ui{reviewed coverage} to hatch the
    cells no reviewed span covers (a cell counts as reviewed when at least half of it lies
    under one). Tick \ui{unreviewed only} to fade the reviewed ones.
  \item \textbf{Select and open a channel.} Click a row to select it; double-click it, or
    press Enter on the selected row, to open it. The bottom bar shows the selected channel's
    counts and two buttons: \ui{Open} \emph{channel} goes to \apppage{explore.signal}, and
    \ui{Cross-channel from} \emph{channel} goes to \apppage{explore.cross-channel}.
\end{enumerate}

\subsection{Where machine and human disagree}\label{sec:explore-disagree}
\ui{colour by disagree} turns the map amber where a run and a person gave different answers
about the same place. Two cases count: \emph{machine yes, human no} (a detection a person
rejected in Review, or one inside reviewed windows marked as not interesting) and
\emph{machine no, human yes} (a span a person said yes to, where a run covered the place and
found nothing). A place no run covered, or no person reviewed, is \emph{not comparable}: it
stays grey and is counted separately in the bottom bar. With no run chosen under
\ui{Detections from}, every run on the recording is pooled, and the map says so.

\screenshot{pages/zzz_x_divergence/explore.corpus--fixup-x-disagree-on-one-run-picked}{Disagree mode
with one run chosen: amber cells are places where that run and a reviewer differ; grey ones
could not be compared.}

\pagefield{Where it leads}
\apppage{explore.signal} for one channel; \apppage{explore.cross-channel} to compare
channels. The recording menu's \ui{Import a recording in Settings › Datasets} link opens
\apppage{settings.datasets}.

\begin{watchout}[A demo marker that is out of date]
With \ui{reviewed coverage} ticked, the caption under the map still carries a \ui{demo} tag.
The hatching is in fact drawn from your database (the reviewed spans the bridge returns), so
the tag is wrong and can be ignored.
\end{watchout}

\begin{underthehood}[Corpus]
Client: \code{webui/client/src/explore/CorpusPage.tsx}, \code{Heatmap.tsx},
\code{RightRail.tsx}, \code{RecordingMenu.tsx}. Reads: \code{GET /api/recordings} and
\code{GET /api/corpus/{file}/coverage} (\code{webui/server/app.py}, which calls
\code{corpus.coverage} in \code{webui/server/corpus.py}); the rail's runs, tags and reviewed
coverage come through \code{getCorpusLive} in \code{webui/client/src/api/explore.ts}, from
\code{GET /api/corpus/{file}/runs} and \code{GET /api/channels/{id}/tags}
(\code{webui/server/explore_routes.py}). The run and method filters are sent as run ids; when
no run matches both, the page sends \code{run=-1} so the detection layers come back empty
rather than unfiltered. The disagreement counts come from
\code{Working/discovery/divergence.py}. The coverage route answers 423 for the held-out file.
\end{underthehood}

\pagesection{explore.signal}{Signal}{One channel at real scale}{live}
Signal shows one channel in three tiers on one time axis: the whole channel, the span you
are looking at, and the one event (the app calls it a \term{motif}{motif}) you have opened.
It is where you choose a span and send it on.

\pagefield{When you'd come here}
After choosing a channel on \apppage{explore.corpus}, or from a channel name anywhere else in
the app. Come here when you want to see what the signal actually looks like, step through
the marked events one at a time, or pick a stretch to analyse.

\screenshot{pages/explore/explore.signal--motif}{A channel with a motif open. The top tier
is the whole channel with the blue span box; the middle tier is that span with its marks;
the bottom tier is the selected motif with its onset and end.}

\pagefield{Doing things here}
\subsection{Find a span}\label{sec:explore-find-span}
\begin{enumerate}
  \item In the top tier, drag the blue box, or drag either of its grips, to choose the
    stretch the middle tier shows. Under the trace, a ribbon shows annotation verdicts by
    colour and another shows how dense the detections are.
  \item In the middle tier, drag to pan and scroll to zoom about the cursor. The keyboard
    does the same: the left and right arrows pan by a tenth, \ui{+} and \ui{-} zoom, \ui{F}
    fits the whole channel. Coloured bands mark annotations (green), detections (blue),
    artifacts (red) and the selected motif (orange).
  \item Choose which runs' detections are drawn with the \ui{detections} chip. Its checklist
    can group runs \ui{by run}, \ui{by method} or \ui{by template}, and has \ui{colour by
    run} and \ui{hide surrogates}.
  \item Click a band to open it in the bottom tier, or walk the marks in time order with the
    arrows beside the counter (keys \ui{[} and \ui{]}). The bottom tier shows the motif with
    $\pm$10, 20 or 60\,s of context (\ui{context}), and an axis in seconds from its onset.
\end{enumerate}

\subsection{Send a span on}\label{sec:explore-send-span}
The row under the tiers, \emph{Selected span}, always describes the stretch in the middle
tier.
\begin{enumerate}
  \item \ui{Send span to Analyse} (key \ui{A}) makes that stretch the source of the chain on
    \apppage{analyse.chain} and takes you there.
  \item \ui{Take span for Review} (key \ui{R}) writes the stretch to the database as a human
    annotation with verdict \emph{seed}, puts it in Review's \emph{Explore spans} queue, and
    offers it as a seed in Discovery. A message with \ui{Open Review} appears.
  \item The bottom tier has the same two acts for the motif alone: \ui{Send motif to
    Analyse} and \ui{Review this motif} (Shift+\ui{R}).
  \item You can add tags (\ui{+ tag}, key \ui{T}) and a note to the span. A tag is lower-case
    words joined by hyphens. A note is never a verdict: verdicts are given in Review.
\end{enumerate}

\subsection{The drawer}\label{sec:explore-drawer}
The four cards along the bottom (\ui{Filters \& search}, \ui{Annotations}, \ui{Detections},
\ui{Keyboard shortcuts}) open a drawer over the lower tiers. It lists every annotation or
detection on the channel in a table you can filter, sort and export with \ui{Export CSV} or
\ui{Export JSON}. Click a row to open it as the motif. Press \ui{D} to open or close the
drawer and \ui{?} for the full keyboard map.

\pagefield{Where it leads}
\apppage{analyse.chain} with the span as its source; \apppage{review.inspector} for a span
taken for Review; \apppage{explore.cross-channel} (key \ui{C}, or the \ui{Signal} /
\ui{Cross-channel} switch in the top bar) to see the same place on other channels.

\begin{watchout}[Some buttons only remember in this browser tab]
\ui{Save span} keeps the span's tags and note in this browser tab only; nothing is written
to the database and a reload loses it. The same is true of the drawer's \ui{Send selected to
Review}, \ui{Add tag} and \ui{Stage for Analyse}: each shows a message, but writes nothing.
\ui{Analyse events} says it is not wired yet. To put a span in front of a reviewer, use
\ui{Take span for Review}, which does write.
\end{watchout}

\begin{watchout}[A recording with no declared unit]
If the recording's file does not say whether it holds volts or millivolts, an amber note
appears and the y axis reads ``?'' instead of mV. The numbers are drawn as stored.
\end{watchout}

\begin{underthehood}[Signal]
Client: \code{webui/client/src/explore/SignalPage.tsx}, with \code{Overview.tsx} (top tier),
\code{SpanView.tsx} (middle), \code{MotifView.tsx} (bottom), \code{SpanActions.tsx},
\code{SignalDrawer.tsx} and \code{DetectionsPicker.tsx}. Reads: \code{GET
/api/channels/{id}}, \code{/window} (a min/max envelope sized to the plot's width) and
\code{/spans} in \code{webui/server/app.py}; runs and drawer rows through \code{getSignalLive}
in \code{webui/client/src/api/explore.ts}. \ui{Take span for Review} posts to
\code{/api/annotations/seed} (\code{take_span_for_review} in
\code{webui/server/explore_routes.py}), which writes through the human-only door
\code{write_human} and finds or creates the open \code{explore-spans} queue. The medoid
overlay and the nearest-family line have no live source yet and stay empty. The page passes
its live read to the header as if it were demo data, so the \ui{demo data} chip shows on
every channel even though the data is yours.
\end{underthehood}

\pagesection{explore.cross-channel}{Cross-channel}{The same window on every channel}{live}
Cross-channel draws one short window on several channels of the same recording, stacked, with
one channel as the reference. For each other channel it gives the time lag and the
correlation with the reference, and sorts the pair into a bin: artifact, propagation or
independent.

\begin{plainterms}[Lag and correlation]
Two microphones in one room hear the same clap a moment apart. The \emph{lag} is that moment:
how far one trace must slide to line up best with the other. The \emph{correlation} $r$ says
how well they line up once slid, from $-1$ (mirror images) through 0 (unrelated) to 1
(identical). If two electrodes show the same shape at the same instant, one wire is probably
picking up the other: an artifact. If the shape arrives a little later on the second
electrode, it may be a signal travelling through the mycelium: propagation.
\end{plainterms}

\begin{plainterms}[Why the page takes the drift out first]
Two people walking up two different hills at the same time are both ``going up'' for an hour,
whatever their steps look like. Two electrodes whose level drifts slowly over ten minutes are
the same: compared as they are, they look almost perfectly related at lag 0, even when their
events have nothing in common. So before measuring the lag and $r$, the page removes from
each channel the straight line that best fits its window (\emph{detrending}). The traces are
still drawn as recorded; only the measurement uses the straightened copy. A straight line
cannot remove a drift that curves, so a curved drift can still look like a shared signal; the
page marks rows where that may be happening.
\end{plainterms}

\pagefield{When you'd come here}
When you have found an event on one channel and want to know whether it is local, appears
elsewhere at the same moment, or travels from one channel to another.

\screenshot{pages/explore/explore.cross-channel--default}{The reference channel on top, with
its span shaded, and the other selected channels below it. Each row ends with its lag, its $r$
and its bin.}

\pagefield{Doing things here}
\begin{enumerate}
  \item \textbf{Choose the window.} The window is one span of the reference channel plus a pad
    each side. The page opens on the channel's most notable human span (a \emph{seed} first,
    then \emph{interesting}, \emph{artifact}, \emph{unsure}; the earliest within each), else
    its first detection, else its first span of any kind; a channel with no span at all opens
    on its first ten minutes. A span longer than ten minutes is cut to
    its first ten. The \ui{window} list offers the default span, up to 40 human spans worth
    anchoring on, and the one in view; the \ui{‹} and \ui{›} arrows step through every span of
    the channel, annotations and detections together, in time order, and \ui{span} \emph{n}
    \ui{of} \emph{N} says where you are. \ui{pad} adds \ui{20 s each side} (the default) or
    \ui{60 s each side}. Above the stack the page says what the window is, for example ``seed
    span at 1.25\,h $\cdot$ 300\,s span + 20\,s each side = 340\,s'', and the time axis counts
    seconds from the span's start.
  \item \textbf{Choose the channels.} \ui{reference} picks the channel the others are
    compared with; \ui{channels} picks up to ten rows. At first the page shows the reference
    and the five other channels with the strongest $|r|$ against it, strongest first, however
    weak. Drag a row's grip, or press Alt and an arrow key, to reorder it. Click a channel
    name to open it in Signal.
  \item \textbf{Read the lags.} Hover to get a crosshair across every row. \ui{align}:
    \ui{lag-aligned} slides each row back by its measured lag so the shapes sit on top of each
    other; a row whose lag is beyond \ui{max lag} stays as recorded and says ``not aligned''.
    Hover a row's $r$ to see the reading as recorded, drift left in, beside the detrended one.
  \item \textbf{Read the bins.} The \emph{Classification for this window} card counts the
    rows in each bin (\emph{artifact}, \emph{propagation}, \emph{independent}, and \emph{no
    match} for a row whose lag and $r$ could not be measured); click a tile to highlight its
    rows. Hover a bin badge for the rule that produced it, as the bridge applied it. A row
    marked \ui{may be curved drift} reads $|r| \ge 0.9$ near lag 0 both detrended and as
    recorded: read its trace before its bin.
  \item \textbf{Choose the y scale.} \ui{y} offers \ui{absolute mV}, \ui{shared gain,
    centred} (the default) and \ui{per channel} (each row scaled to fit, marked
    \ui{normalised}). In the centred view every row has the same millivolts per row height:
    the widest row's own range, which fills its row and says ``sets the gain''. Each row is
    drawn about the middle of its own range, so none is clipped and depths still compare.
    Electrodes can sit volts apart in baseline while their events are a few millivolts high,
    so a shared absolute scale draws every trace flat. Every row prints its own true range at
    the right.
\end{enumerate}

\paragraph{A window opened from a link.} A link with an exact start and end
(\texttt{?t0=} and \texttt{?t1=}) shows that window unpadded, as a button you click to return to
the channel's own spans. When the link came from an edge on \apppage{library.family}, the
window is measured as recorded, without detrending, so the page reproduces the lag, $r$ and bin
the edge stores; the note above the stack then reads ``lag $\cdot$ r as recorded''. Any other
exact window is detrended like a picked one, so the same window gives the same answer whichever
link opened it.

\pagefield{Where it leads}
\ui{Apply a template across these} \emph{N} \ui{channels} opens Discovery with the channels
filled in (\apppage{discovery.runs}). \ui{Classify every} \emph{family} \ui{member in
Library} runs only when you came here from a family's edge on \apppage{library.family}; it
starts a job that classifies every member of that family across channels. Opened any other
way, the card is greyed and prints why under it.

\begin{watchout}[A bin here is a lead, not a finding]
The bins on this page come from the lag and $r$ alone. The chance test and the noise-floor
test that the family classification from Library runs are not run on a window, and each bin's
tooltip now says so. With the default settings a pair is an artifact when
$|{\rm lag}| \le 1$\,s and $|r| \ge 0.98$, propagation when $1 < |{\rm lag}| \le 50$\,s and
$|r| \ge 0.5$, and independent otherwise. The lag is searched over the whole window, so
unrelated channels can land in \emph{propagation} by chance: the project's own simulation (the
note for research question 6) put about one pair in six of independent random walks there. A
straight line does not remove a curved drift; the same simulation found curved drifts binned
\emph{artifact} almost every time after a straight-line detrend, which is what the \ui{may be
curved drift} mark is for. The bridge can also take out a cubic instead of a line
(\texttt{detrend=cubic} on \texttt{/api/cross}), but the page does not offer it yet.
\end{watchout}

\begin{watchout}[Controls that do less than they look]
\ui{max lag} only changes which rows lag-aligned moves; it does not recompute anything, and the
progress bar that appears when you change it is a short animation. The \ui{lag} list's
\ui{computed on whole channel} option is not available yet, and nor is the \ui{display}
list's detrended view. \ui{Review this window on all channels} only records the request in
this browser tab. The \ui{shared ground} chip in the channel picker is drawn for one fixed
pair of channel names, not read from Settings.
\end{watchout}

\begin{underthehood}[Cross-channel]
Client: \code{webui/client/src/explore/CrossChannelPage.tsx}; the window, the detrend choice,
the centred gain and the lag shift are decided in \code{crossModel.ts}
(\code{resolveCrossWindow}, \code{crossDetrendFor}, \code{centredStack}, \code{lagShift}), the
y geometry in \code{crossScale.ts}. Read: \code{getCrossChannel} in
\code{webui/client/src/api/explore.ts} resolves the window and calls \code{getCrossWindow},
\code{GET /api/cross/{id}} with \texttt{t0}, \texttt{t1} and \texttt{detrend} (\code{get_cross_channel} in
\code{webui/server/explore_routes.py}). The bridge cuts the same absolute window from every
sibling channel, strides it down past 200{,}000 samples, and makes one call per pair to
\code{Working.cross_channel.classify_waveforms} with the page's detrend (\texttt{linear}, or
\texttt{none} for an edge's window). That call removes each window's trend
(\code{_remove_trend}), takes the peak of the z-normalised cross-correlation and bins it with
\code{bin_for} under the rule from Settings › Analysis defaults (\code{rule_from_settings}),
with its chance and floor arguments left at their pass-through defaults. When it detrends, a
second call as recorded rides along as \code{as_recorded}, never binning, and
\code{drift_suspect} sets the curved-drift mark. A channel that is nothing but its trend raises
\code{FlatWindow} and reads \emph{no match} with the reason. The words behind each bin are
\code{CrossChannelRule.describe_window}. Held-out siblings are skipped. A row whose data cannot
be read says so on the row instead of being drawn as a second reference.
\end{underthehood}

\pagesection{explore.span-edit}{Span edit}{Editing a span, opened from Review}{preview}
Span edit is where you will redraw the start and end of a span a machine found, without
changing what the machine found. It is reached from Review's \ui{Edit span in Explore}
button (key \ui{E}) and from \ui{Redraw in Explore} on a Library family. Today the page runs
only on one example span from the demo data and saves nothing.

\pagefield{When you'd come here}
When a candidate in Review is the right event but its edges are wrong: it starts too early,
or cuts off the tail.

\screenshot{pages/explore/explore.span-edit--edited}{The original extent in grey behind your
new extent in blue. The card below lists what saving would write: a new revision, not a
change to the detection.}

\pagefield{Doing things here}
\begin{enumerate}
  \item Drag either blue grip, or type a sample number into \ui{start} or \ui{end}. The
    small arrows and the arrow keys nudge by one sample (Shift for ten).
  \item \ui{snap to} chooses where a dragged edge lands: \ui{steepest sample},
    \ui{trough}, \ui{zero crossing} or \ui{free}.
  \item Read \emph{What saving writes}: a new revision of the span, kept beside the
    machine's original, in the human store.
  \item \ui{Save and return to Review}, or \ui{Return without saving}.
\end{enumerate}

\pagefield{Where it leads}
Back to \apppage{review.inspector}.

\begin{watchout}[Preview only]
Everything on this page comes from \code{webui/client/src/fixtures/explore.ts}. Only the
span \code{m-1846} exists; opening any other span, which is what Review's button does for a
real candidate, shows ``No span \ldots{} in the demo data''. Saving stays in this browser
tab. The design it shows is the intended one: a person's redrawn edges go to the annotation
store as a new revision, and the detection row is never touched.
\end{watchout}

\begin{underthehood}[Span edit]
Client: \code{webui/client/src/explore/SpanEditPage.tsx}; data from \code{getSpanEdit} in
\code{webui/client/src/api/explore.ts}, which returns \code{spanEditDemo} from the fixtures.
The core already has revisions (\code{Working/library/revisions.py}, \code{add_revision});
nothing connects this page to it yet.
\end{underthehood}
